The hall emptied in two ways and he had stopped being able to keep them separate.

There had been four hundred and some at the worst of it, laid in rows the length of the Great Hall with lamps burning low and the fallen gallery beam standing at the eastern end like a thing in a cave. By the time the smoke off Rhaenys’s Hill had gone from green to grey to nothing there were a hundred and sixty. That was a good number. Sam knew it was a good number; he could have told anybody, with the Citadel’s own figures to hand, that a hundred and sixty out of four hundred and some after a fortnight of that was better than any account he had ever read of any siege in the last two hundred years.

The difficulty was that the number had come down two ways at once. Some of them had got up and walked out and gone back to whatever was left of their houses. Some of them had died. And after the first week Sam had found that he could no longer reliably remember which of the empty places had been which, and he had begun keeping a separate sheet for it, and then he had begun to be frightened of the separate sheet.

He kept it anyway. He kept everything.

\sceneBreak

The register was in the small hall behind the kitchens on a trestle and it was his and he had not asked for it.

That had happened the way most things happened: Ser Roland’s men had read four hundred and eleven names at the head of the first pit and somebody had had to write them, and Sam had been writing everything since the burial council, so Sam wrote them, and after that there was a book with the names of the dead of King’s Landing in it in his hand, and after that people came to look at it.

They came all day. They came in a queue that formed itself before dawn and never entirely dispersed, and the queue was patient and orderly, and every single person in it wanted the same thing, which was to be told a name was in the book.

He was extremely bad at it.

He understood this. He had understood it by about the fourth morning and had gone to Lord Tyrion and asked to be relieved of it and had been told, kindly and immediately, that there was nobody, and had gone back to the trestle.

The trouble was that the book was not what people thought it was. The book was a record of names spoken aloud at the head of a pit by men who had been given those names by other men who had received a body from a person who had said a name. There were four hundred and eleven in the first reading and Sam could have told you, of those, perhaps three hundred and thirty where a family member had handed over a body and given a name and the name had been written at the arch by a clerk who could read. The other eighty were names that had come to him along a chain of two or three people, in the cold, at speed, and about eleven of them he privately believed were the same name written twice with different spellings.

That was what the book was. And a woman would come to the trestle and say a name and Sam would find it, or not find it, and either way she would want to know what it meant, and the truthful answer was that it meant a name had been spoken, and that was all it meant, and nobody in that queue had come four miles to hear that.

\sceneBreak

The woman came in the middle of the fish, and she was about forty, and her name was Marra, and she had been in the queue since before light for the fourth time.

“Hesk,” she said. “Hesk, of Pisswater Bend, a caulker. He’s got a mark on his jaw here.”

“I’ve looked for you three times.”

“Look again.”

He looked again. He looked properly, the whole of both readings, and he was fast at it by then and it took perhaps four minutes, and there was no Hesk and no Hesken and no Esk, and there was a Hersk that he had already established with her on the second morning was a bargeman of sixty from Cobbler’s Square whose sister had come in person.

“He isn’t there.”

“Then put him there.”

Sam looked up.

“I’m sorry?”

“Put him in.” She said it perfectly reasonably. She had thought about it; he could see that she had thought about it for some days. “You’ve the pen. He went out of the house on the second night and he’s not come back and there’s nothing in this city he’d have gone to instead of coming back, and I’ve had three weeks of it, and my girls ask me. So put him in your book and I’ll take them up the hill and I’ll show them where he is and that’ll be the end of it.”

“I can’t do that.”

“Why?”

And this was the moment — he thought about it a very great deal afterward, and always came back to this precise moment, with the pen in his hand and the queue going out of the door behind her — at which he had four or five things available to him and chose the worst one.

“Because it wouldn’t be true,” Sam said.

\sceneBreak

He did not say it unkindly. He wanted that noted somewhere and there was nowhere to note it. He said it the way he would have said it to a novice at the Citadel who had asked whether one might reasonably round a figure, and he went on, because he had begun and because he had spent his entire life believing that if you explained a thing properly to somebody they would see it.

“This book is the only thing in this city that says who died. In ten years it will be the only thing. There’ll be men who come back — there are men coming back now, off ships, out of the Reach, and they’ll walk into this hall and want to know, and this is what they’ll be given, and if there are names in it that were put there because somebody asked me to then it isn’t a record any more, it’s a — “ He heard himself and could not stop himself. “It’s a list of things people needed. Which is a different document. And I can make you that document, I’ll make it now, I’ll write his name on a separate sheet and I’ll keep it with this one and I’ll write on it exactly what it is, and anybody who ever looks at the one will be told the other exists, and that’s — “

“He’s dead,” Marra said.

“I know.”

“You \textit{know} he’s dead.”

“I think he’s almost certainly dead and I have no way of establishing it and that is not the same — “

The queue had gone quiet behind her. He became aware of it about a heartbeat too late.

“My lord Maester,” Marra said, and she was not shouting, and that was the part that got into him and stayed. “There’s a hundred people behind me and I’ve been here four mornings and I’ve watched you do this to six of them. You’re the only man in this city who’s told me the truth about anything and I’ve come back to you three times because of it. And you’re sat there with a pen telling me my husband isn’t dead enough for your book.”

“That is not what I — “

“It’s what it is.”

\sceneBreak

It was Bett’s cousin who fixed it, in about nine seconds, and Sam did not even see her come in.

She was perhaps nineteen and she carried food and water in that wing and she had spent a fortnight watching him and she came round the end of the trestle, took a spare sheet off his own stack, said “Give me the pen, ser,” took it out of his hand before he had decided anything, and wrote at the top of it, in a hand considerably better than his:

\textbf{MISSING SINCE THE SECOND NIGHT. NOT FOUND. NOT IN THE PITS THAT ANYBODY KNOWS OF.}

Then she said: “Hesk. Caulker. Pisswater Bend, mark on the jaw,” and wrote it, and turned the sheet round.

“That’s on the table beside the book,” she said to Marra. “Every day. Same table. And if any man comes off a ship in ten years asking after your Hesk, he’ll be shown both, because Maester Tarly writes everything down and he never loses anything, and that’s the truth of it and I’ve watched him do it for a fortnight.” She put the pen back in Sam’s hand. “That do?”

Marra of Pisswater Bend looked at the sheet for a while.

“That’ll do,” she said.

And she went, and the queue moved, and Sam sat there with the pen in his hand and understood that the girl had given that woman precisely what he had been trying to give her, in nine seconds, with no argument, having conceded nothing whatsoever, and that the entire difference between them was that she had written the second sheet \textit{first}.

\textit{I had it,} he thought. \textit{I had the answer. It was the third thing I said to her. And I put it fourth in the sentence, after ‘it wouldn’t be true’, and by then she couldn’t hear it.}

He was not much good in a room. He had known that since he was nine years old and had a very great deal of evidence for it, and he had somehow persuaded himself over the past three weeks that he had stopped being that, on the grounds that he had got better at telling one person on a pallet a difficult thing.

He wrote the second sheet’s heading into the record that night with the date on it, and noted, in his own account, whose idea it had been. It seemed the least he could do and he was aware it was not much.

\sceneBreak

He could not write to Gilly.

He had tried five times and the difficulty was not what he had expected it to be. He had expected it to be difficult in the ordinary way — the news, the amount of it, the fact that Jon’s hand and Bran and the arch and four hundred and eleven names could not be got into a letter of any length that a woman alone in White Harbor with a small boy ought to have to read on her own.

That was not the difficulty. The difficulty was that he did not know whether she was alive.

Nothing had come south of the Neck in three weeks. Not a bird, not a ship, not a rider. Lady Sansa had four out and none back and had stopped saying so out loud in front of people. And Sam had sat down five times with a good sheet and a sharpened pen, and each time had got as far as \textit{My love} and then had understood, in a way that arrived new every single time and never once became familiar, that he was writing to somebody about whom he possessed exactly no information.

The fifth time he got a great deal further. He wrote for about an hour. He told her about the girl of fifteen with the hands, whose name was Wenda, and who had held a retractor for four hours on the first night and had not once looked away, and who he intended to have properly trained if there was ever a Citadel again. He told her that he had cried in front of Jon Snow and that Jon had never mentioned it. He told her a thing about Bran that he had not told anybody.

Then he read it back and burned it, because there was a fair chance that the only person who would ever read that letter was a stranger going through a dead woman’s box in White Harbor, and he found he could not have that.

He had sent the other kind of letter as well — not to her, but for her, addressed to whoever kept the lists of the sheltered at White Harbor, asking only for two names and whether they could be found among them, nothing more, no explanation, because an explanation would have cost more paper than the asking was worth. That one he had not burned. It had gone out on a northbound ship a fortnight past, and there was nothing further to be done to it now except wait on an answer that had not yet come, and might, for all he knew, arrive built from the same silence as everything else out of that country.

\textit{Lord Davos goes north after the thaw,} he thought. \textit{Nine weeks. He’ll go up that coast and put in at White Harbor and he’ll come back and tell me, and until then there is not one single thing I can do about it, and I have four hundred names to copy.}

He went and copied names. It worked, mostly. That was the discovery of that whole month and he was not proud of it: that work did in fact function as a place to put a thing, that it functioned extremely well, and that this was almost certainly the mechanism by which his father had lived to be fifty-eight.

\sceneBreak

He found the dagger’s entry by accident, looking for something else entirely, and it was four lines long.

He had gone into the lower armoury of the Red Keep to find out about spears. That was all. Ser Roland wanted two hundred men equipped for the pits and the water lines and nobody knew what was in the building, and there were three chests of ledgers in a dry room off the armoury that had not been opened since Robert Baratheon’s time and possibly rather longer, and Sam had gone down there for an afternoon with a lamp because he was the only person in King’s Landing who would voluntarily do such a thing.

The spear ledgers were there and were useless. Underneath them was an inventory of the royal armoury made in the reign of Aegon the Fifth, in a beautiful hand, by a clerk who had cross-referenced everything.

Sam read it for two hours and did not once think about the dagger.

Then, on the eleventh folio, under \textbf{VALYRIAN STEEL, THE PROPERTY OF THE HOUSE}, among nine entries of which seven were swords:

\textit{Item, a dagger, the blade of Valyrian steel, the hilt of dragonbone, no maker’s mark, of no great age. Not of the old inheritance. Bought or received in the reign of Aegon III and entered under his hand, the purpose not stated.}

And in the margin, in a different and later ink, in the tight hand of some clerk perhaps a hundred years after:

\textit{Not found at the inventory of 259. Sought. See the third book, which is lost.}

\sceneBreak

Sam sat in a dry room off the lower armoury of the Red Keep with a lamp going down and read those four lines eleven or twelve times.

He had the dagger. It was in his satchel and had been since the second day, when he had taken it out of Jon Snow’s fist with a pair of forceps and cleaned what was underneath it and had then, without asking anybody’s permission and without any clear reason he could have stated, not given it back.

He got it out and put it on the ledger.

Dragonbone hilt. No maker’s mark. Not of great age — and the clerk had been right about that, whatever \textit{great age} meant to a man cataloguing Valyrian steel; the thing had no history on it at all, no name, no scabbard worth the making, which was exactly the sort of blade a man buys when he wants a good blade and not a famous one.

\textit{Entered under his hand. The purpose not stated.}

\textit{Not found at the inventory of two fifty-nine.}

Sam Tarly was not a fanciful man and had spent a great deal of his life being told he was, and what he did next was what he always did, which was to write down precisely what he had and precisely what he did not.

What he had: that a Valyrian dagger with a dragonbone hilt and no maker’s mark had been royal Targaryen property, entered personally by a Targaryen king; that it had gone out of the Red Keep’s inventory at some point before the year 259; that the record of how it went was in a book that was lost; and that a dagger answering that description had been used in an attempt on Brandon Stark’s life at Winterfell some seven years past, and had thereafter been given to Arya Stark, and had been dropped into a wildfire trench below the Mud Gate, and had been picked up out of the fire by Jon Snow.

What he did not have: anything at all connecting the one to the other.

He wrote both lists. He wrote the second list longer than the first, because it was longer.

Then he sat there with a dragonbone hilt under his hand in a cold room and thought about a thing he had read out to a queen and a bastard in a solar at Winterfell — \textit{Rhaegar and Lyanna, wed at Starfall, the child born at the tower and given the name Aegon} — and about the fact that this dagger had been bought by a Targaryen king for a purpose nobody had written down, and had come out of that armoury by a route recorded in a lost book, and had gone into a boy’s chest in a bedchamber in the North.

He did not follow it any further than that. He wanted it on the record afterward that he had not followed it, and it was on the record, because he wrote that down too.

There were a hundred and sixty people in the Great Hall and Wenda had been on her own with them for four hours and it was long past when he should have gone up.

He put the dagger back in the satchel and took the folio with him.

\sceneBreak

He went to the Tower of the Hand last, as he did every night, which was two hundred and eleven steps and got no easier.

Lady Sansa had done the evening work already; he could tell from the room. The cloth was folded on the sill in the particular way she folded it and the body had been turned onto its left side and the jug was empty.

He did what he had come for, which took twenty minutes: the chest, the pulse, the skin over the hip and the shoulder where the pressure was, the mouth. There was nothing new. There had been nothing new for three weeks, which he had at first found encouraging and had latterly come to understand differently.

Fourteen breaths to the minute. It had been fourteen since the first night.

“I found something,” Sam said. “It’s not — I don’t know what it is. It might be nothing at all.”

He got the folio out and read the four lines aloud, and then the marginal note, and then sat back on his heels.

There had been a period of about a week when he had done this every night and had waited afterward, without admitting to himself that he was waiting.

He no longer waited. That was the change and he had not been able to identify when it happened. He came up two hundred and eleven steps every night and did the work and told him things, and it had become, at some point in the fish, a thing he did for himself.

\textit{Which is all right,} Sam thought. \textit{That’s allowed. People do that at graves and nobody thinks anything of it.}

He put the folio away.

“The trouble is that I’ve nobody to ask,” he said. “That’s the whole of it, if I’m honest. There’s four things in my satchel that oughtn’t to be in one man’s satchel and there’s nobody in this city I can put any of them in front of, and you were — “

He stopped.

Fourteen to the minute, steady, in the dark, at the top of the Tower of the Hand.

“You’d have been unbearable about the ledger,” Sam said. “You’d have known which book it was.”

He put the lamp out and went down the two hundred and eleven steps and back to the hall, where a hundred and sixty people were waiting, and got on with it, because there was a great deal of it and it was his and there was nobody else.